% =====================================================================
% DRAFT SECTION FOR PAPER 47 -- NOT INTEGRATED INTO main.tex
% Editor draft, 2026-10-13. Owner decides.
% Sources (phase 186, only what survived hostile review with label PROVED):
%   R186_042 Thm 2.1, Thm 2.2, Thm 2.3 (Halmos; sign slip corrected here), Prop 2.5;
%   R186_043 Lemma M0 (Krein formula; R186_115: rank-one value is a LOWER bound of beta_1,
%            kappa of the slope law is fitted -> only the PROVED part enters), Cor M0-step;
%   R186_044 Thm 1 (C_T necessary), Thm 2 (exact Schur), Lemma Q, Prop K; Lemma P only ARGUED;
%   R186_045 Lemma T, Lemma T' (moments: margin <= k(1-beta_1));
%   R186_047 Lemmas 1-3, Theorem (L-W); R186_116: values "PROVED modulo convergent floating
%            quadrature", consequence only in the three windows;
%   R186_048 Lemma 3.1 (form unchanged at T_n), Lemma 2.3 (newborn components, Shannon < 1/2);
%   R186_049 Prop 2.1 (old/new partition);
%   R186_028 Thm 1.2/1.4, Thm 2.1 (index by sectors; R186_106/107 labels).
% Every statement was checked against the paper's own definitions: eq:symbol, eq:form
% (form domain D_T), eq:W, P_T, thm:A/eq:Aform, lem:K, lem:mbounds, lem:contA (and the
% paragraph before it: H_T, compact resolvent), lem:monoT, prop:NC (S_{T,tau_1}),
% prop:band, prop:negtrace (Pos_T, Neg_T), thm:W, lem:roughEF, thm:parityschur,
% prop:polarcap, thm:log5A, thm:log5W, thm:margin, def:W/def:PD (not used, recorded).
% See CHECK_TABLE.md, part IV. Lines "% EDITOR FLAG" must be resolved before
% integration. Compiles standalone (preamble skipped when \input into main.tex).
% Cross-references to main.tex and to the other draft sections (sec:inertia, sec:coer)
% are undefined in the standalone build, as in the previous drafts.
% =====================================================================
\ifdefined\PT\else
\documentclass[11pt]{amsart}
\usepackage{amsmath,amssymb,amsthm,mathtools,mathrsfs}
\usepackage[margin=1.08in]{geometry}
\usepackage{booktabs}
\usepackage[hypertexnames=false]{hyperref}
\hypersetup{hidelinks}
\newtheorem{theorem}{Theorem}[section]
\newtheorem{proposition}[theorem]{Proposition}
\newtheorem{lemma}[theorem]{Lemma}
\newtheorem{corollary}[theorem]{Corollary}
\theoremstyle{definition}
\newtheorem{definition}[theorem]{Definition}
\theoremstyle{remark}
\newtheorem{remark}[theorem]{Remark}
\newcommand{\C}{\mathbb C}
\newcommand{\R}{\mathbb R}
\newcommand{\Rel}{\operatorname{Re}}
\newcommand{\Iml}{\operatorname{Im}}
\newcommand{\supp}{\operatorname{supp}}
\newcommand{\PT}{\mathcal P_T}
\newcommand{\IT}{I_T}
\newcommand{\lamA}{\lambda_{\min}(A_T|\PT)}
\newcommand{\lamW}{\lambda^{\mathrm W}_T}
\newcommand{\sT}{\sigma_T}
\newcommand{\Fh}{\widehat F}
\newcommand{\cF}{c_F}
\newcommand{\NT}{\mathcal N_T}
\newcommand{\Aconst}{\mathsf A_T}
\newcommand{\PiT}{\Pi_T}
\newcommand{\DraftStandalone}{}
\begin{document}
\title{Draft section for Paper 47: the window form as a truncated Wiener--Hopf operator}
\maketitle
\fi

\newcommand{\Pos}{\mathrm{Pos}_T}
\newcommand{\Neg}{\mathrm{Neg}_T}
\newcommand{\fp}{\mathfrak p}
\newcommand{\fn}{\mathfrak n}
\newcommand{\fa}{\mathfrak a_T}
\newcommand{\DT}{\mathcal D_T}
\newcommand{\tr}{\operatorname{tr}}

\section{The window form as a truncated Wiener--Hopf operator: inertia crossings, prolate reduction and moment bounds}\label{sec:wh}

This section records the structure of $\fa$ on $\PT$ as the difference of its positive and
negative parts (Proposition~\ref{prop:negtrace}), read as a one-parameter family of truncated
convolution operators whose inertia changes at the eigenvalues of a trace-class pencil. Every
statement is a reformulation or a reduction; none of them carries a sign. Where a reduction
produces a finite object (a matrix of size $2\cdot\#$components, a Schur complement, a list of
moments), that object inherits the margin of Theorem~\ref{thm:margin} in full
(Lemmas~\ref{lem:wh-T}, \ref{lem:wh-Tprime} and Remark~\ref{rem:wh-honest}). The word
``Wiener--Hopf'' refers only to the shape of the operators $R_Tf(D)R_T$; no Wiener--Hopf
index is used or defined (the symbol $\sT-\lambda$ is not elliptic, Theorem~\ref{thm:W}).

\emph{Notation.} $\DT=\{F\in L^2(\IT):\int\log(2+|\tau|)|\Fh|^2<\infty\}$ is the form domain of
\eqref{eq:form}; $R_T$ is multiplication by $\mathbf 1_{\IT}$ on $L^2(\R)$ and $f(D)$ the Fourier
multiplier $\widehat{f(D)F}=f\Fh$, so that $\langle R_Tf(D)R_TF,F\rangle=\frac1{2\pi}\int f|\Fh|^2$
for $F\in L^2(\IT)$; in particular $\fa(F)=\langle R_T\sT(D)R_TF,F\rangle$. As in
Proposition~\ref{prop:negtrace}, $\sT=\sigma_+-\sigma_-$ with $\sigma_\pm=\max(\pm\sT,0)$,
$\Pos=R_T\sigma_+(D)R_T$ (a closed form on $\DT$) and $\Neg=R_T\sigma_-(D)R_T$ (a nonnegative
trace-class operator on $L^2(\IT)$). For a measurable $E\subset\R$, $S_E:=R_T\mathbf 1_E(D)R_T$ is the
time--band limiting operator, $0\le S_E\le1$, with $\langle S_EF,F\rangle=\frac1{2\pi}\int_E|\Fh|^2$;
for $E=[-\tau_1,\tau_1]$ this is $S_{T,\tau_1}$ of Proposition~\ref{prop:NC}, and $\tr S_E=T|E|/\pi$
when $|E|<\infty$ (kernel $\frac1{2\pi}\int_Ee^{i\tau(t-s)}d\tau$ on the diagonal). By
Proposition~\ref{prop:band}, the open set $U_T:=\{\sT<0\}\subset\R$ is a finite union of bounded
open intervals, symmetric under $\tau\mapsto-\tau$; its connected components are
$E_0=(-a_0,a_0)$ and $E_i=J_i\cup(-J_i)$, $1\le i\le n_T-1$, with $J_i\subset(0,\infty)$ disjoint
intervals; we write $n_T$ for the number of components of $U_T\cap[0,\infty)$ and $s_i:=\max_{E_i}\sigma_-$.
$\Pi$ is the orthogonal projection of $L^2(\IT)$ onto $\PT$, $JF(t)=F(-t)$ the reflection
(it preserves $\PT$ and commutes with every $R_Tf(D)R_T$ with $f$ even, cf.\ the paragraph before
Lemma~\ref{lem:contA}), and $h^*_T$ a unit minimizer of $\fa$ on $\PT\cap\DT$ (attained: $H_T$ has
compact resolvent, same paragraph). $\lambda_{\mathrm A}(T):=\lamA$.

\subsection{The pencil and its inertia crossings}\label{subsec:wh-pencil}

\begin{lemma}[the positive part is bounded below]\label{lem:wh-cP}
For every $\delta>0$ put $c_\delta:=\min\{\sT(\tau):|\tau|\ge\tau_{\mathrm{last}}+\delta\}>0$ and
$\lambda_0(\delta):=\|S_{T,\tau_{\mathrm{last}}+\delta}\|<1$. Then for all $F\in\DT$,
\[
\Pos(F)\ \ge\ c_\delta\big(1-\lambda_0(\delta)\big)\|F\|^2 .
\]
Hence $c_P:=\inf\{\Pos(F):F\in\PT\cap\DT,\|F\|=1\}>0$, and the closed form $\Pos|_{\PT\cap\DT}$ has an
associated self-adjoint operator $\fp\ge c_P$ on $\PT$, with $\fp^{1/2}$ a bijection of $\PT\cap\DT$ onto $\PT$.
\end{lemma}
\begin{proof}
$\sT>0$ on $|\tau|>\tau_{\mathrm{last}}$ (Proposition~\ref{prop:band}(iv)), $\sT$ is continuous and
$\sT\ge m_{\mathrm{lo}}-\Aconst\to+\infty$ (Lemma~\ref{lem:mbounds}(ii),(iii)), so the minimum
$c_\delta$ over the closed set $|\tau|\ge\tau_{\mathrm{last}}+\delta$ exists and is positive. Then
$\sigma_+\ge c_\delta\mathbf 1_{|\tau|\ge\tau_{\mathrm{last}}+\delta}$ pointwise, so
$\Pos(F)\ge c_\delta\frac1{2\pi}\int_{|\tau|\ge\tau_{\mathrm{last}}+\delta}|\Fh|^2
=c_\delta\big(\|F\|^2-\langle S_{T,\tau_{\mathrm{last}}+\delta}F,F\rangle\big)$, and $\lambda_0(\delta)<1$
is shown in the proof of Proposition~\ref{prop:negtrace}(v) (an eigenvector with eigenvalue $1$ would
have $\Fh=0$ on a set of positive measure, impossible for an entire $\Fh\ne0$). Closedness of
$\Pos$ on $\DT$: $0\le\sigma_+\le|C_0|+\log(1+|\tau|)+\Aconst$ (Lemma~\ref{lem:mbounds}, as in the
proof of Lemma~\ref{lem:roughEF}) and $\sigma_+\ge\frac12\log(2+|\tau|)$ for $|\tau|$ large
(Lemma~\ref{lem:mbounds}(ii)), so $\Pos+\|\cdot\|^2$ is equivalent to the weighted norm
$\int\log(2+|\tau|)|\Fh|^2$, exactly as for $\fa$ in the paragraph before Lemma~\ref{lem:contA}. The
last claim is the spectral theorem for $\fp\ge c_P>0$.
\end{proof}

\begin{definition}\label{def:wh-B}
Let $\fn:=\Pi\Neg\Pi|_{\PT}$ (bounded, $\fn\ge0$, trace class, $\tr\fn\le\tr\Neg=\frac{2T}{2\pi}\int\sigma_-$
by Proposition~\ref{prop:negtrace}(i)), and
\[
B_T:=\fp^{-1/2}\fn\,\fp^{-1/2}\quad\text{on }\PT,\qquad \beta_1\ge\beta_2\ge\dots\ (\downarrow0)\ \text{its eigenvalues},
\]
so that $B_T\ge0$ is trace class with $\tr B_T\le\tr\fn/c_P$. For $\beta>0$ let
$\sigma_\beta:=\sigma_+-\sigma_-/\beta$ and let $W_T(\sigma_\beta)$ denote the form $\Pos-\Neg/\beta$ on
$\PT\cap\DT$ (the truncated convolution operator with symbol $\sigma_\beta$, compressed to $\PT$);
$W_T(\sigma_1)=\fa|_{\PT}$. For a closed lower-bounded form $\mathfrak q$ on a subspace $V\cap\DT$,
$n_-(\mathfrak q|V)$ is the largest dimension of a subspace on which $\mathfrak q$ is negative definite
(Lemma~\ref{lem:parity}).
\end{definition}

\begin{theorem}[inertia crossings]\label{thm:wh-pencil}
\emph{(i)} For $u\in\PT\cap\DT$ and $\beta>0$,
\[
W_T(\sigma_\beta)(u)=\big\langle(1-\beta^{-1}B_T)\fp^{1/2}u,\fp^{1/2}u\big\rangle .
\]
Hence $n_-(W_T(\sigma_\beta)|\PT)=\#\{j:\beta_j>\beta\}<\infty$, the map $\beta\mapsto n_-(W_T(\sigma_\beta)|\PT)$
is nonincreasing, and it jumps exactly at the $\beta_j$ (by the multiplicity of $\beta_j$). In particular
\[
n_-(\fa|\PT)=\#\{j:\beta_j>1\},\qquad \lambda_{\mathrm A}(T)>0\iff\beta_1<1,\qquad \lambda_{\mathrm A}(T)\ge0\iff\beta_1\le1 .
\]
\emph{(ii)} $1-\beta_1=\inf\{\fa(u)/\Pos(u):u\in\PT\cap\DT,\ u\ne0\}=\varepsilon^*/(1+\varepsilon^*)$, where
$\varepsilon^*:=\sup\{\varepsilon\ge0:\fa\ge\varepsilon\,\fn\text{ on }\PT\}$; and
$\beta_1=\sup\Neg(u)/\Pos(u)$ is attained at $u_1:=\fp^{-1/2}v_1/\|\fp^{-1/2}v_1\|$, $v_1$ a unit eigenvector of $\beta_1$.
\emph{(iii)} If $\lambda_{\mathrm A}(T)\ge0$ then, with $u_1$ and $h^*_T$ unit vectors,
\begin{equation}\label{eq:wh-sandwich}
\frac{\lambda_{\mathrm A}(T)}{\Pos(u_1)}\ \le\ 1-\beta_1\ \le\ \frac{\lambda_{\mathrm A}(T)}{\Pos(h^*_T)}\ \le\ \frac{\lambda_{\mathrm A}(T)}{c_P}.
\end{equation}
\emph{(iv)} $n_-(\fa|\PT)\le\#\{j:\nu_j(\fn)>c_P\}$, where $\nu_j(\fn)$ are the eigenvalues of $\fn$.
\end{theorem}
\begin{proof}
(i) For $u\in\PT\cap\DT$, $\Pos(u)=\|\fp^{1/2}u\|^2$ and $\Neg(u)=\langle\fn u,u\rangle=\langle B_T\fp^{1/2}u,\fp^{1/2}u\rangle$,
which gives the identity. Since $\fp^{1/2}$ maps $\PT\cap\DT$ onto $\PT$ (Lemma~\ref{lem:wh-cP}), a subspace on
which $W_T(\sigma_\beta)$ is negative definite corresponds to one on which $1-\beta^{-1}B_T$ is negative
definite and conversely; for the compact $B_T$ the latter has largest dimension $\#\{j:\beta_j>\beta\}$
(min--max). The last three statements are the case $\beta=1$: if $\beta_1<1$ then
$\fa(u)\ge(1-\beta_1)\|\fp^{1/2}u\|^2\ge(1-\beta_1)c_P\|u\|^2$; if $\beta_1\ge1$ then $\fa(\fp^{-1/2}v_1)=(1-\beta_1)\le0$;
similarly for $\ge0$. (ii) $\fa(u)/\Pos(u)=1-\Neg(u)/\Pos(u)$ and $\sup\Neg/\Pos=\sup_v\langle B_Tv,v\rangle/\|v\|^2=\beta_1$,
attained at $v_1$. $\fa\ge\varepsilon\fn\iff\fp\ge(1+\varepsilon)\fn\iff B_T\le(1+\varepsilon)^{-1}\iff\beta_1\le(1+\varepsilon)^{-1}$,
so $\varepsilon^*=\beta_1^{-1}-1$ and $\varepsilon^*/(1+\varepsilon^*)=1-\beta_1$. (iii) $\fa(u_1)=(1-\beta_1)\Pos(u_1)\ge\lambda_{\mathrm A}$;
and $\lambda_{\mathrm A}=\fa(h^*)=\Pos(h^*)-\Neg(h^*)\ge(1-\beta_1)\Pos(h^*)$ since $\Neg\le\beta_1\Pos$; finally $\Pos(h^*)\ge c_P$.
(iv) On $\PT\cap\DT$, $\fa(u)\ge c_P\|u\|^2-\langle\fn u,u\rangle$; if $\fa<0$ on a subspace of dimension $n$,
then $\fn>c_P$ there, and by min--max $\nu_n(\fn)>c_P$.
\end{proof}
% EDITOR FLAG (R186_034 Thm 2.1(iv)/Thm 2.2): the memo writes "lambda_A <= (1-beta_max) Pos(h*)" and the
% interval [lambda_A/Pos(h*), lambda_A/c_P] for 1-beta_max. The two inequalities that follow from the
% definitions are (eq:wh-sandwich): the lower end carries the PENCIL extremal u_1, the upper end the FORM
% minimizer h*; both ends are <= lambda_A/c_P. Numerically u_1 and h* coincide to the digits printed, so
% no number changes. R186_049 Prop 2.1 has the correct orientation.

\begin{remark}\label{rem:wh-pencil}
(1) Theorem~\ref{thm:wh-pencil} is Sylvester's law of inertia written for the pair $(\Pos,\Neg)$ of
Proposition~\ref{prop:negtrace}(iv); it adds to that proposition the identification of the negative
index of $\fa$ with the number of pencil eigenvalues above $1$, and of the bottom with the last inertia
crossing of the family $\beta\mapsto W_T(\sigma_\beta)$. The question ``$\lambda_{\mathrm A}(T)>0$'' is
the question ``is the last crossing $\beta_1$ below $1$'', an extremal-eigenvalue question and not a
counting one; counting theorems with error $o(\log T)$ (Landau--Widom type) locate none of the
$\beta_j$ to the precision needed (R186\_039, R186\_042; CITED, not used).
(2) Part (iv) improves the counts of Proposition~\ref{prop:NC} at the computed windows
($4$ at $\frac12\log3$, $10$ at $\frac12\log5$ against $37$, $54$; NUMERICAL, R186\_034), but it is a bound on
the index, not on the bottom.
(3) \emph{Inheritance.} By \eqref{eq:wh-sandwich}, $1-\beta_1$ is $\lambda_{\mathrm A}(T)$ divided by a quantity
in $[c_P,\Pos(u_1)]$; at $T=\frac12\log5$ this is $1.7\cdot10^{-11}$ against $\lambda_{\mathrm A}=3.9\cdot10^{-12}$
and $\Pos(h^*)=0.22$ (NUMERICAL, R186\_034/042/043; the certified value is Theorem~\ref{thm:log5A}).
\end{remark}

\begin{proposition}[the pencil in the eigenbasis of the negative part]\label{prop:wh-basis}
Let $(\varphi_k,\nu_k)_{k\ge1}$ be an orthonormal eigenbasis of $\Neg$ on $L^2(\IT)$ with $\nu_k>0$
(its range closure), and $\varphi_k$ for $k\le0$ an orthonormal basis of $\ker\Neg$. Every $\varphi_k$ with
$\nu_k>0$ lies in $\DT$ and $P_{kl}:=\Pos(\varphi_k,\varphi_l)=\frac1{2\pi}\int\sigma_+\widehat{\varphi_k}\overline{\widehat{\varphi_l}}$
is finite. For $u=\sum_ka_k\varphi_k\in\PT\cap\DT$,
\[
\fa(u)=\sum_{k,l}P_{kl}a_k\bar a_l-\sum_{k\ge1}\nu_k|a_k|^2,\qquad
1-\beta_1=\inf_{u\in\PT\cap\DT\setminus0}\frac{\sum_{k,l}P_{kl}a_k\bar a_l-\sum_{k\ge1}\nu_k|a_k|^2}{\sum_{k,l}P_{kl}a_k\bar a_l}.
\]
\end{proposition}
\begin{proof}
$\Neg\varphi=\nu\varphi$ with $\nu>0$ gives $\varphi=\nu^{-1}R_T\sigma_-(D)R_T\varphi$, the restriction to $\IT$ of the
entire function $(\sigma_-\widehat\varphi)^\vee$ (bounded compactly supported multiplier), hence smooth on $\overline{\IT}$,
so $\widehat\varphi(\tau)=O(1/|\tau|)$ and $\varphi\in\DT$. The two displays are the definition of $\fa=\Pos-\Neg$ and
Theorem~\ref{thm:wh-pencil}(ii).
\end{proof}
This is the form in which R186\_042 Thm 2.2 verified the pencil numerically (the matrix $P$ of
``weighted prolates'' of Widom; $1-\beta_1$ agrees with the direct computation to $10^{-4}$ relative at
$T=\frac12\log5$; NUMERICAL). It is a change of basis and nothing more; the $\varphi_k$ with $\nu_k=0$ are
needed because $\PT\cap\DT$ is not spanned by the range of $\Neg$.

\subsection{Constant weights: principal angles, and why they do not decide}\label{subsec:wh-angles}

\begin{theorem}[two bands with constant weights]\label{thm:wh-angles}
Let $\Omega_-,\Omega_+\subset\R$ be disjoint measurable even sets with $|\Omega_-\cup\Omega_+|>0$,
$s_\pm>0$, $Q_\pm:=\mathbf 1_{\Omega_\pm}(D)$, $Q:=Q_-+Q_+$, and consider on $L^2(\IT)$ the forms
$\mathfrak n(u):=s_-\langle R_TQ_-R_Tu,u\rangle$, $\mathfrak p(u):=s_+\langle R_TQ_+R_Tu,u\rangle$. Put
$V:=Q\,L^2(\IT)\subset\operatorname{ran}Q$ and, for $v\in V\setminus0$, let $\phi(v)\in[0,\frac\pi2]$ be the angle
between $v$ and $\operatorname{ran}Q_-$, $\cos^2\phi(v)=\|Q_-v\|^2/\|v\|^2$. Then $u\mapsto Qu$ is injective on $L^2(\IT)$,
and for every $u\ne0$ with $Q_+u\ne0$,
\[
\frac{\mathfrak n(u)}{\mathfrak p(u)}=\frac{s_-}{s_+}\cot^2\phi(Qu).
\]
Consequently, on any subspace $X\subset L^2(\IT)$: $\sup_X\mathfrak n/\mathfrak p=\frac{s_-}{s_+}\cot^2\phi_{\min}(QX)$ with
$\phi_{\min}$ the smallest angle between $QX$ and $\operatorname{ran}Q_-$, and
$n_-\big((\mathfrak p-\mathfrak n)|X\big)=\max\{\dim Y:Y\subset QX,\ \|Q_-v\|^2>\theta\|v\|^2\ \forall v\in Y\setminus0\}$,
$\theta:=s_+/(s_++s_-)$.
\end{theorem}
\begin{proof}
If $Qu=0$ then $\widehat u=0$ a.e.\ on $\Omega_-\cup\Omega_+$, a set of positive measure, so the entire function
$\widehat u$ vanishes and $u=0$. For $v=Qu$: $Q_-v=Q_-u$ and $Q_+v=Q_+u$ ($Q_\pm Q=Q_\pm$), and
$v=Q_-v+Q_+v$ orthogonally, so $\|Q_+v\|^2=\sin^2\phi(v)\|v\|^2$. Thus
$\mathfrak n(u)/\mathfrak p(u)=s_-\|Q_-v\|^2/(s_+\|Q_+v\|^2)$. The negativity condition
$s_+\|Q_+v\|^2<s_-\|Q_-v\|^2$ is $(s_++s_-)\|Q_-v\|^2>s_+\|v\|^2$.
\end{proof}
% EDITOR FLAG (R186_042 Thm 2.3): the memo states beta_j = (s_-/s_+) tan^2 phi_j with phi_j the angles
% between V and ran Q_-. The decomposition Q_- - beta Q_+ = (1/2)[(1-beta)Q + (1+beta)J] has the sign +(1+beta)J
% (the memo wrote -(1+beta)J), and the correct value is cot^2 of the angle to ran Q_- (= tan^2 of the angle to
% ran Q_+). The memo's own Corollary 2.4 (beta_j = (s_-/s_+) lambda_j/(1-lambda_j)) is the cot^2 form and is
% correct; the two statements of the memo are inconsistent with each other by a complementary angle.

\begin{corollary}[no gap: the two parts commute]\label{cor:wh-nogap}
If $\Omega_+=\R\setminus\Omega_-$ with $|\Omega_-|<\infty$, then $Q=1$, $V=L^2(\IT)$, $R_TQ_-R_T=S_{\Omega_-}$ and
$R_TQ_+R_T=1-S_{\Omega_-}$ commute, and with $\lambda_0\ge\lambda_1\ge\dots$ the eigenvalues of $S_{\Omega_-}$ the pencil
$(\mathfrak n,\mathfrak p)$ on $L^2(\IT)$ has eigenvalues $\beta_j=\frac{s_-}{s_+}\frac{\lambda_j}{1-\lambda_j}$, so that
$n_-(\mathfrak p-\mathfrak n|L^2(\IT))=\#\{j:\lambda_j>\theta\}$ and $\mathfrak p-\mathfrak n>0$ on $L^2(\IT)$ iff
$s_-\lambda_0<s_+(1-\lambda_0)$. On $\PT$ (codimension $2$) the pencil eigenvalues interlace:
$\beta_j(L^2)\ge\beta_j(\PT)\ge\beta_{j+2}(L^2)$.
\end{corollary}
\begin{proof}
$\cos^2\phi(v)$ ranges over the Rayleigh quotients of the compact operator $S_{\Omega_-}$ on $L^2(\IT)$; the
interlacing is min--max for $B$ restricted to a subspace of codimension $2$.
\end{proof}

\begin{proposition}[the bottom is not a function of the band, the negative profile and the angles]\label{prop:wh-notangles}
For $\varepsilon>0$ let $\sigma^{(\varepsilon)}:=\varepsilon\sigma_+-\sigma_-$. Then $\{\sigma^{(\varepsilon)}<0\}=U_T$,
$\sigma^{(\varepsilon)}_-=\sigma_-$, the operators $S_{U_T}$ and all principal angles of the pair $(R_T,\mathbf 1_{U_T}(D))$
are unchanged, and the top pencil eigenvalue of $\sigma^{(\varepsilon)}$ on $\PT$ is $\beta_1/\varepsilon$; hence the
form of $\sigma^{(\varepsilon)}$ on $\PT$ is indefinite for $\varepsilon<\beta_1$ and positive for $\varepsilon>\beta_1$.
\end{proposition}
\begin{proof}
$\Pos^{(\varepsilon)}=\varepsilon\Pos$ and $\Neg^{(\varepsilon)}=\Neg$, so $B^{(\varepsilon)}_T=B_T/\varepsilon$.
\end{proof}
This is the content of R186\_042 Prop.~2.5 in its simplest form: the sign is carried by the relation
between $\sigma_+$ and $\sigma_-$ (through the whole of both profiles), not by the geometry of
$U_T$ alone. The companion fact that constant-weight and level-per-component models of $\sT$ with the
true $U_T$ give $1-\beta_1$ of either sign and of size $10^{-2}$--$1$ against $10^{-11}$ for $\sT$ itself
is NUMERICAL (R186\_042 \S3.3) and is not entered.

\subsection{The Kre\u{\i}n formula for a rank-one negative part}\label{subsec:wh-krein}

\begin{lemma}[Kre\u{\i}n]\label{lem:wh-krein}
Let $\fp\ge c>0$ be a self-adjoint operator on a Hilbert space $H$ with form $\fp(u)=\|\fp^{1/2}u\|^2$, let $\fn\ge0$ be
bounded with $\fn\psi=\nu\psi$, $\|\psi\|=1$, $\nu>0$, and let $\beta_1:=\|\fp^{-1/2}\fn\fp^{-1/2}\|$.
\emph{(i)} $\beta_1\ge\nu\langle\fp^{-1}\psi,\psi\rangle$, with equality if $\fn=\nu\,\psi\otimes\psi$.
\emph{(ii)}
\[
\langle\fp^{-1}\psi,\psi\rangle=\sup_{u\ne0}\frac{|\langle u,\psi\rangle|^2}{\fp(u)}=\Big(\inf_{u\perp\psi}\fp(\psi+u)\Big)^{-1}.
\]
\emph{(iii)} If $\fn=\nu\,\psi\otimes\psi$, then
$\fp-\fn>0$ on $H$ $\iff$ $\nu\langle\fp^{-1}\psi,\psi\rangle<1$ $\iff$ $\nu<\inf_{u\perp\psi}\fp(\psi+u)$.
\end{lemma}
\begin{proof}
(i) $\fn\ge\nu\,\psi\otimes\psi$ (both sides are $\ge0$ and agree on $\psi$; $\fn-\nu\psi\otimes\psi$ is $\fn$ on
$\psi^\perp$), so $\fp^{-1/2}\fn\fp^{-1/2}\ge\nu\,\fp^{-1/2}\psi\otimes\fp^{-1/2}\psi$, whose norm is $\nu\|\fp^{-1/2}\psi\|^2$;
equality in the rank-one case. (ii) Cauchy--Schwarz in the inner product $\langle\fp\cdot,\cdot\rangle$:
$|\langle u,\psi\rangle|^2=|\langle\fp^{1/2}u,\fp^{-1/2}\psi\rangle|^2\le\fp(u)\langle\fp^{-1}\psi,\psi\rangle$, with equality at
$u=\fp^{-1}\psi$; and $\{w:\langle w,\psi\rangle=1\}=\psi+\psi^\perp$, so $\inf_{u\perp\psi}\fp(\psi+u)=\inf_{\langle w,\psi\rangle=1}\fp(w)
=1/\sup|\langle w,\psi\rangle|^2/\fp(w)$. (iii) $\fp-\nu\psi\otimes\psi=\fp^{1/2}(1-\nu\fp^{-1/2}\psi\otimes\fp^{-1/2}\psi)\fp^{1/2}$,
and a rank-one perturbation $1-\nu\,e\otimes e$ is positive iff $\nu\|e\|^2<1$; the second equivalence is (ii).
\end{proof}
% EDITOR FLAG (R186_115 §1.1): R186_043 presents nu<p^{-1}psi,psi> as "deciding the sign" of the model; for a
% negative part of rank > 1 it is a LOWER bound of beta_1 (part (i)), so "nu<p^{-1}psi,psi> < 1" is a
% necessary condition only. Entered with that orientation.

\begin{corollary}[one negative interval on a constant background]\label{cor:wh-step}
Let $E=(-W,W)$, $a,b_0>0$, $\sigma=b_0\mathbf 1_{\R\setminus E}-a\mathbf 1_E$, and $\lambda_0\ge\lambda_1\ge\dots$ the eigenvalues
of $S_E=S_{T,W}$. Then $R_T\sigma(D)R_T=b_0(1-S_E)-aS_E$ on $L^2(\IT)$, its pencil eigenvalues are
$\beta_j=\frac a{b_0}\frac{\lambda_j}{1-\lambda_j}$, and the form is positive on $L^2(\IT)$ iff $a\lambda_0<b_0(1-\lambda_0)$.
More generally, if $\sigma=b\mathbf 1_{\R\setminus E}-a\mathbf 1_E$ with $b\ge b_0$ on $\R\setminus E$ and $\psi$ a unit top eigenvector of $S_E$, then
\[
b_0(1-\lambda_0)\ \le\ \inf_{u\perp\psi}\ \langle R_Tb(D)R_T(\psi+u),\psi+u\rangle\ \le\ \frac1{2\pi}\int_{\R\setminus E}b\,|\widehat\psi|^2 .
\]
\end{corollary}
\begin{proof}
The first part is Corollary~\ref{cor:wh-nogap}. For the second, $R_Tb(D)R_T\ge b_0(1-S_E)$ (order of multipliers),
so Lemma~\ref{lem:wh-krein}(ii) gives $\langle(R_Tb(D)R_T)^{-1}\psi,\psi\rangle\le(b_0(1-\lambda_0))^{-1}$, which is the left
inequality; the right one is $u=0$.
\end{proof}

\begin{remark}\label{rem:wh-model}
(1) In the model of Corollary~\ref{cor:wh-step} the sign is decided by one number per component,
$a\lambda_0(S_E)$ against $b_0(1-\lambda_0(S_E))$. For $\sT$ this mechanism is absent: $\sigma_+$ vanishes at
$\partial U_T$ (linearly), so $b_0=0$ in every component, and level-per-component minorants of $\sigma_+$ give
pencils that are indefinite (NUMERICAL, R186\_042 \S3.3, R186\_043 \S7). The scaling law
``$\beta_1\approx\kappa\,a\,w\,T^2/b_1$'' for a background $b_1\operatorname{dist}(\tau,E)$ is NUMERICAL and ARGUED
(R186\_043 \S3), with $\kappa$ fitted and varying between $0.55$ and $1.24$ over the parameters
(R186\_115 \S1.2); it is not entered.
(2) The components of $U_T$ are sub-Shannon at every computed window: $T|J_i|/\pi\le0.95$, and the time--band
operator $S_{U_T}$ has no eigenvalue near $1$ ($\lambda_0(S_{U_T})\in[0.76,0.88]$ for $T\in[0.62,1]$; NUMERICAL,
R186\_042 \S3.2). So $1-\beta_1\approx10^{-11}$ is not a plateau value of any time--band operator; the
asymptotics of Slepian--Fuchs--Widom and Landau--Widom are outside their hypotheses here (CITED, R186\_042 \S1)
and are not used anywhere in this section.
\end{remark}

\subsection{Reduction to the component prolates and the exact Schur complement}\label{subsec:wh-prolate}

\begin{definition}\label{def:wh-V}
For each component $E_i$ of $U_T$ and each parity $\pm$ let $\lambda_0^{(i),\pm}\ge\lambda_1^{(i),\pm}\ge\dots$ be the
eigenvalues of $S_{E_i}$ on the even ($+$) and odd ($-$) subspaces of $L^2(\IT)$, and $\psi_i^\pm$ a unit eigenvector
of $\lambda_0^{(i),\pm}$ (the index-$0$ prolate of the component in its parity sector; $S_{E_i}$ commutes with $J$
since $E_i$ is even). Put
\[
V_T:=\operatorname{span}\{\psi_i^\pm\},\qquad V_T^P:=\Pi V_T\subset\PT,\qquad W_T:=\PT\ominus V_T^P,
\]
$H_T:=\fa|_{V_T^P}$ (a real symmetric form on a space of dimension $\le2n_T$), $A_{WW}:=\fa|_{W_T\cap\DT}$ and
$A_{VW}$ the mixed part. $C_T$ is the statement $H_T\succ0$. Variant: $V_T':=V_T+\operatorname{span}\{\varphi_i^\pm\}$,
$\varphi_i^\pm$ a unit top eigenvector of $\mathrm{Neg}_{T,i}:=R_T\sigma_-\mathbf 1_{E_i}(D)R_T$ in each parity sector, and
$V_T'^P,W_T'$ accordingly.
\end{definition}
% EDITOR FLAG (R186_115 MAJOR-2): R186_043 wrote the condition on V cap P_T and computed it on Pi V; the two
% differ (the literal intersection has dimension 2n_T-2 and does not detect the falsifier M_1 at log 2).
% R186_044 adopted V^P = Pi V, which is what enters here.

\begin{lemma}\label{lem:wh-W}
\emph{(i)} $V_T\subset\DT$, $V_T^P\subset\DT$, and $W_T=\PT\cap V_T^\perp$. \emph{(ii)} $V_T$, $V_T^P$ and $W_T$ are
$J$-invariant, so $\fa|_{W_T}$ splits into even and odd parts.
\end{lemma}
\begin{proof}
(i) $\psi=\lambda_0^{-1}S_{E_i}\psi$ is the restriction of an entire function, so $\psi\in\DT$ as in
Proposition~\ref{prop:wh-basis}; $\Pi\psi=\psi-$(Tate component), and $e^{\pm t/2}\mathbf 1_{\IT}\in\DT$. For $F\in\PT$,
$\langle F,\Pi v\rangle=\langle\Pi F,v\rangle=\langle F,v\rangle$. (ii) The $\psi_i^\pm$ are $J$-eigenvectors, $J$ preserves $\PT$ and commutes with $\Pi$.
\end{proof}

\begin{theorem}[the prolate block is necessary]\label{thm:wh-necessary}
$\lambda_{\mathrm A}(T)\le\lambda_{\min}(H_T)$, $n_-(H_T)\le n_-(\fa|\PT)$, and $\sup_{V_T^P}\Neg/\Pos\le\beta_1$. In particular
$C_T$ is necessary for $\lambda_{\mathrm A}(T)>0$, and if $H_T$ has a negative eigenvalue then $\lambda_{\mathrm A}(T)<0$.
\end{theorem}
\begin{proof}
Rayleigh quotients on the subspace $V_T^P\subset\PT\cap\DT$; Cauchy interlacing for the compression; Theorem~\ref{thm:wh-pencil}(ii).
\end{proof}

\begin{lemma}[the negative part on $W_T$]\label{lem:wh-neg}
For $F\in W_T$,
\[
\Neg(F)\ \le\ \mathrm{Neg}_{\mathrm{ub}}\,\|F\|^2,\qquad
\mathrm{Neg}_{\mathrm{ub}}:=\min\Big\{\max_\pm\sum_is_i\lambda_1^{(i),\pm},\ \Big\|\sum_is_i\widetilde S_{E_i}\Big\|\Big\},
\]
where $\widetilde S_{E_i}:=S_{E_i}-\sum_\pm\lambda_0^{(i),\pm}\,\psi_i^\pm\otimes\psi_i^\pm\ge0$.
For $F\in W_T'$ the same holds with $s_i\lambda_1^{(i),\pm}$ replaced by $\nu_1^{(i),\pm}$, the second eigenvalue of $\mathrm{Neg}_{T,i}$ in the sector.
\end{lemma}
\begin{proof}
$\sigma_-=\sum_i\sigma_-\mathbf 1_{E_i}$ (disjoint supports), so $\Neg=\sum_i\mathrm{Neg}_{T,i}$, and $\sigma_-\mathbf 1_{E_i}\le s_i\mathbf 1_{E_i}$
gives $\mathrm{Neg}_{T,i}\le s_iS_{E_i}$. Write $F=F_++F_-$ by parity; $F_\pm\in W_T$ (Lemma~\ref{lem:wh-W}) and
$\langle S_{E_i}F,F\rangle=\sum_\pm\langle S_{E_i}F_\pm,F_\pm\rangle$. In the sector $\pm$, $F_\pm\perp\psi_i^\pm$ and $S_{E_i}$ is compact
self-adjoint, so $\langle S_{E_i}F_\pm,F_\pm\rangle\le\lambda_1^{(i),\pm}\|F_\pm\|^2$ (spectral theorem); summing over $i$ and $\pm$
gives the first bound. Also $\langle S_{E_i}F,F\rangle=\langle\widetilde S_{E_i}F,F\rangle$ for $F\perp V_T$, and $\widetilde S_{E_i}\ge0$
because the $\psi_i^\pm$ are orthonormal eigenvectors of $S_{E_i}$; the second bound is the norm of the sum. The variant is the
same argument with $\mathrm{Neg}_{T,i}$ in place of $s_iS_{E_i}$.
\end{proof}

\begin{lemma}[the positive part on $W_T$ by sublevel layers]\label{lem:wh-pos}
For $s\ge0$ let $L_s:=\{\tau:\sT(\tau)\le s\}$ (bounded, a finite union of intervals, $L_0=\overline{U_T}$, nondecreasing in $s$) and
\[
\widetilde S_{L_s}:=S_{L_s}-\sum_i\sum_\pm\lambda_0^{(i),\pm}\,\psi_i^\pm\otimes\psi_i^\pm,\qquad \bar\mu(s):=\|\widetilde S_{L_s}\| .
\]
Then $\widetilde S_{L_s}\ge0$ and $s\mapsto\bar\mu(s)$ is nondecreasing, and for every $F\in W_T\cap\DT$, $\|F\|=1$, and every
partition $0=s_0<s_1<\dots<s_N$,
\[
\Pos(F)=\int_0^\infty\big(1-\mu_F(L_s)\big)\,ds\ \ge\ \sum_{k=0}^{N-1}(s_{k+1}-s_k)\big(1-\bar\mu(s_{k+1})\big)^+=:\mathrm{Pos}_{\mathrm{lb}},
\]
where $\mu_F(A):=\frac1{2\pi}\int_A|\Fh|^2$.
\end{lemma}
\begin{proof}
$L_s$ is bounded by Proposition~\ref{prop:band}(i) and Lemma~\ref{lem:mbounds}(ii) ($\sT\to+\infty$) and has finitely many
components by real-analyticity of $\sT$ (Proposition~\ref{prop:band}(iii)). For $\|F\|=1$, $\mu_F$ is a probability measure
and $\Pos(F)=\int\sigma_+d\mu_F=\int_0^\infty\mu_F(\{\sigma_+>s\})ds$ (layer-cake), with $\{\sigma_+>s\}=\R\setminus L_s$ for
$s\ge0$. All $E_i\subset L_0\subset L_s$, so $S_{L_s}\ge\sum_iS_{E_i}\ge\sum_i\sum_\pm\lambda_0^{(i),\pm}\psi_i^\pm\otimes\psi_i^\pm$,
which gives $\widetilde S_{L_s}\ge0$; the subtracted term does not depend on $s$, so $\widetilde S_{L_s}$ and $\bar\mu$ are nondecreasing.
For $F\in W_T$, $\mu_F(L_s)=\langle S_{L_s}F,F\rangle=\langle\widetilde S_{L_s}F,F\rangle\le\bar\mu(s)$, and for $s\in(s_k,s_{k+1}]$,
$1-\mu_F(L_s)\ge1-\mu_F(L_{s_{k+1}})\ge(1-\bar\mu(s_{k+1}))^+$; integrate.
\end{proof}
% EDITOR FLAG (R186_116 MINOR-1): R186_047 claims "mu-bar(0) = sup_{V^perp} mu_F(U)" (exactness at s=0). False:
% the psi of different components are not orthogonal, so S-tilde_U differs from (1-Pi_V)S_U(1-Pi_V) (15 % at log 2).
% The lemma is a valid bound, not an identity; nothing here claims exactness.

\begin{theorem}[coercivity on the prolate complement]\label{thm:wh-LW}
For every $T>0$, with the constants of Lemmas~\ref{lem:wh-neg} and \ref{lem:wh-pos},
\[
\fa(F)\ \ge\ c_W\,\|F\|^2\quad(F\in W_T\cap\DT),\qquad c_W:=\mathrm{Pos}_{\mathrm{lb}}-\mathrm{Neg}_{\mathrm{ub}} ,
\]
and likewise on $W_T'$ with the variant constant. Both constants are spectral data of time--band operators of
$(T,U_T,\sT)$: no minimizer and no cancellation enters, and the same bound holds for every symbol of the class
(any $\sigma=\sigma_+-\sigma_-$ with $\{\sigma<0\}$ a finite union of bounded intervals).
\end{theorem}
\begin{proof}
$\fa=\Pos-\Neg$ and the two lemmas.
\end{proof}

\begin{theorem}[the exact Schur complement over the prolate block]\label{thm:wh-schur}
Suppose $A_{WW}\ge c\|\cdot\|^2$ on $W_T\cap\DT$ for some $c>0$ (Theorem~\ref{thm:wh-LW} whenever $c_W>0$).
\emph{(i)} For each $e\in V_T^P$ there is a unique $Ke\in W_T\cap\DT$ with $\fa(Ke,w)=\fa(e,w)$ for all $w\in W_T\cap\DT$, and
$\fa(w+e)=\fa(w+Ke)+\mathfrak s_T(e)$, $\mathfrak s_T(e):=\fa(e)-\fa(Ke)=H_T(e)-\mathfrak c_T(e)$, where
$\mathfrak c_T(e):=\sup_{w\in W_T\cap\DT}[2\Rel\fa(e,w)-\fa(w)]\ge0$ is the real quadratic form $A_{VW}A_{WW}^{-1}A_{WV}$ on $V_T^P$.
\emph{(ii)} $\lambda_{\mathrm A}(T)>0\iff\mathfrak s_T\succ0$, and $\lambda_{\mathrm A}(T)\ge0\iff\mathfrak s_T\succeq0$. If $H_T\succ0$,
put $\rho_T:=\max\{\mathfrak c_T(e)/H_T(e):e\in V_T^P\setminus0\}$; then
\[
\lambda_{\mathrm A}(T)>0\iff\big[A_{WW}\ge c>0\big]\ \wedge\ \big[H_T\succ0\big]\ \wedge\ \big[\rho_T<1\big],
\]
where the first bracket is automatic when $\lambda_{\mathrm A}(T)>0$ (with $c=\lambda_{\mathrm A}(T)$).
\emph{(iii)} With $\lambda_S:=\lambda_{\min}(\mathfrak s_T)$ and $\|K\|$ the norm of $K:V_T^P\to W_T$,
$\frac{\min\{c,\lambda_S\}}{(1+\|K\|)^2}\le\lambda_{\mathrm A}(T)\le\min\{c,\lambda_S\}\le\lambda_{\min}(H_T)$, and, if
$\lambda_{\mathrm A}(T)>0$ and $\lambda_S\le c$,
\[
1-\rho_T\ \le\ \frac{\lambda_S}{\lambda_{\min}(H_T)}\ \le\ \frac{(1+\|K\|)^2}{\lambda_{\min}(H_T)}\,\lambda_{\mathrm A}(T).
\]
\end{theorem}
\begin{proof}
(i)--(iii) are Theorem~\ref{thm:coer-schur} of \S\ref{sec:coer} with $(\mathcal V,E)=(W_T\cap\DT,V_T^P)$: $V_T^P$ is
finite-dimensional and contained in $\DT$ (Lemma~\ref{lem:wh-W}), $\PT\cap\DT=(W_T\cap\DT)\oplus V_T^P$, and $\fa$ is an
inner product on $W_T\cap\DT$ equivalent to the weighted norm (Lemma~\ref{lem:mbounds}, as there), so $K$ exists by Riesz.
The equivalence in (ii) is Haynsworth's inertia formula (with $A_{WW}\succ0$, $\fa|\PT\succ0$ iff the Schur complement
$\mathfrak s_T=H_T-\mathfrak c_T\succ0$) followed by Lemma~\ref{lem:wh-krein} for the pencil $(\mathfrak c_T,H_T)$ on the
finite-dimensional $V_T^P$: $H_T-\mathfrak c_T\succ0$ iff $H_T\succ0$ and $\rho_T<1$. For the last display, let $e_S$ be a unit
eigenvector of $\lambda_S$: $1-\rho_T=\min_e\mathfrak s_T(e)/H_T(e)\le\lambda_S/H_T(e_S)\le\lambda_S/\lambda_{\min}(H_T)$, and
$\lambda_S\le(1+\|K\|)^2\lambda_{\mathrm A}(T)$ by the sandwich.
\end{proof}

\begin{remark}[numbers, labels, and what the reduction does not do]\label{rem:wh-LWnum}
(1) \emph{Constants (R186\_047, hostile R186\_116).} Galerkin values: $c_W=+0.41,+0.34,+0.49$ at
$T=\frac12\log3,\log2,\frac12\log5$ (against $\lambda_{\min}(A_{WW})=0.79,1.37,1.14$); with $V_T'$: $+0.55,+0.40,+0.53$
(NUMERICAL). ``Certified'' values, in which $V_T$ is replaced by the explicit Galerkin prolates (Legendre
polynomials), every quantity is an explicit integral, and no inequality ``Galerkin $\le$ true'' is used:
$c_W=+0.136,+0.201$ at $\frac12\log3,\log2$ and $+0.594,+0.548,+0.318$ with $V_T'$ at the three windows.
Their label is PROVED modulo convergent floating-point quadrature ($\ge6$ digits, refinement by a factor $4$;
no interval arithmetic), R186\_116 MAJOR-1; the value $+0.004$ at $\frac12\log5$ with $V_T$ is not used.
(2) \emph{Scope of the equivalence.} Theorem~\ref{thm:wh-schur}(ii) in the form
``$\lambda_{\mathrm A}(T)>0\iff C_T\wedge\rho_T<1$'' with the hypothesis $A_{WW}\ge c>0$ discharged is established only at
$T\in\{\frac12\log3,\log2,\frac12\log5\}$ (with $V_T'$ at the last), where $c_W>0$ has been evaluated; $c_W(T)>0$ is not
proved for all $T$ ($\sum_is_i\lambda_1^{(i)}$ grows with the number of components while $\mathrm{Pos}_{\mathrm{lb}}$ does not).
(3) \emph{Inheritance.} $1-\rho_T$ inherits $\lambda_{\mathrm A}(T)$ by (iii): $1-\rho_T=0.21,\ 3.7\cdot10^{-4},\ 1.0\cdot10^{-8}$
at the three windows (NUMERICAL, R186\_044). $\lambda_{\min}(H_T)=4.3\cdot10^{-3},5.9\cdot10^{-5},5.8\cdot10^{-7}$ has a margin of its
own but is blind at the scale of $\lambda_{\mathrm A}$: on the modified symbols of R186\_114 (a quartet of zeros at
$\frac12\pm\delta\pm i\gamma_1$ added to the zero side, cf.\ Lemma~\ref{lem:wh-Q}) the bottom changes sign at
$\delta_A\approx2.5\cdot10^{-3}$ while $H_T$ stays positive up to $\delta_V\approx3.4\cdot10^{-2}$ at $T=\log2$ (NUMERICAL,
R186\_115 \S3.4, R186\_044 \S4). Adding the index-$1$ prolates to $V_T$ lowers $\lambda_{\min}(H_T)$ to $\lambda_{\mathrm A}$ within
$1$--$3\%$ (NUMERICAL, R186\_115 \S3.3): there is no rank reduction. The correction $\mathfrak c_T$ is not small:
$\|A_{VW}\|\approx0.6$, $\|\mathfrak c_T\|\approx0.16$--$0.19$ against $\lambda_{\min}(H_T)$ (NUMERICAL), so no bound by norms
replaces (ii). None of this asserts a sign.
\end{remark}

\subsection{A quartet on the zero side as a rank-$(2,2)$ form}\label{subsec:wh-quartet}

\begin{lemma}\label{lem:wh-Q}
For real $F\in L^2(\IT)$, $z=\delta+i\gamma$ with $\delta,\gamma\in\R$, and $\Phi_F(s)=\int F(t)e^{(s-\frac12)t}dt$ as in
\eqref{eq:W}, put $Q_z(F):=4\Rel\big[\Phi_F(\tfrac12+z)\Phi_F(\tfrac12-z)\big]$ (the contribution of the four points
$\frac12\pm\delta\pm i\gamma$ to $\sum_\rho\Phi_F(\rho)\Phi_F(1-\rho)$ when they are added to the zero set). Then
\[
Q_z(F)=4|\Phi_c(F)|^2-4|\Phi_s(F)|^2,\qquad \Phi_c(F):=\int F(t)\cosh(\delta t)e^{i\gamma t}dt,
\]
\[
\Phi_s(F):=\int F(t)\sinh(\delta t)e^{i\gamma t}dt ,
\]
a real quadratic form of inertia $(2,2)$ on real $L^2(\IT)$ when $\delta\gamma\ne0$; on even $F$ it is $4\Phi_c(F)^2+4\,\Iml\Phi_s(F)^2$
with $\Phi_c$ real and $\Phi_s$ imaginary, on odd $F$ the roles are exchanged, and $\Phi_s(F)=\delta\int F(t)\,t\,e^{i\gamma t}dt+O(\delta^3)$.
\end{lemma}
\begin{proof}
For real $F$, $\Phi_F(\frac12-z)=\overline{X_-}$ with $X_\pm:=\int Fe^{\pm\delta t}e^{i\gamma t}=\Phi_c\pm\Phi_s$, and
$\Rel[X_+\overline{X_-}]=|\Phi_c|^2-|\Phi_s|^2+\Rel[\Phi_s\overline{\Phi_c}-\Phi_c\overline{\Phi_s}]$, the last term being purely imaginary.
The four real functionals $\Rel\Phi_c,\Iml\Phi_c,\Rel\Phi_s,\Iml\Phi_s$ are linearly independent (the functions
$\cosh\delta t\cos\gamma t$, $\cosh\delta t\sin\gamma t$, $\sinh\delta t\cos\gamma t$, $\sinh\delta t\sin\gamma t$ are). Parity and the
expansion of $\sinh$ are immediate.
\end{proof}
This is Lemma~\ref{thm:inertia-quartet} of \S\ref{sec:inertia} in closed form (there: $4\Rel A_z(F_{\mathrm e})^2-4\Rel A_z(F_{\mathrm o})^2$).

\begin{proposition}[Kre\u{\i}n threshold for a quartet on a finite block]\label{prop:wh-K}
Let $X\subset\PT\cap\DT$ be finite-dimensional with $\fa|_X\succ0$, and let $S_X:X\to\R^2$, $S_Xe=(\Rel\Phi_s(e),\Iml\Phi_s(e))$,
$H_+:=\fa|_X+4|\Phi_c|^2\succ0$. Then
\[
(\fa+Q_z)|_X\succ0\iff 4\,S_XH_+^{-1}S_X^{\mathsf T}\prec I_2\iff \max\operatorname{spec}\big(4S_XH_+^{-1}S_X^{\mathsf T}\big)<1 .
\]
\end{proposition}
\begin{proof}
$(\fa+Q_z)|_X=H_+-4S_X^{\mathsf T}S_X$ by Lemma~\ref{lem:wh-Q}, and for $H_+\succ0$ and $G:=2S_X$, $H_+-G^{\mathsf T}G\succ0$ iff
$\|H_+^{-1/2}G^{\mathsf T}\|<1$ iff $GH_+^{-1}G^{\mathsf T}\prec I$ (Lemma~\ref{lem:wh-krein} for a perturbation of rank $\le2$).
\end{proof}
The first-order law $\lambda_{\mathrm A}\mapsto\lambda_{\mathrm A}+4|\widehat{h^*}(\gamma)|^2-4\delta^2|\int h^*te^{i\gamma t}|^2+O(\delta^4)$
of R186\_044 (Lemma P) is ARGUED (perturbation of a simple eigenvalue, second order not bounded with constants) and
NUMERICAL ($1$--$2\%$); it is not entered.

\subsection{Moments of the pencil and the margin of any finite criterion}\label{subsec:wh-moments}

\begin{lemma}[moment criterion]\label{lem:wh-mom}
Let $m_k:=\tr B_T^k=\sum_j\beta_j^k$, $k\ge1$ (finite, $B_T$ trace class). Then $\beta_1\le m_k^{1/k}$ for all $k$, $m_k^{1/k}\downarrow\beta_1$,
and $\beta_1<1\iff m_k<1$ for some $k$.
\end{lemma}
\begin{proof}
$\beta_1^k\le m_k\le\beta_1^k\sum_j(\beta_j/\beta_1)^k$ and the sum decreases to the multiplicity of $\beta_1$; if $\beta_1<1$ then
$m_k\to0$ by dominated convergence ($\sum_j\beta_j<\infty$).
\end{proof}

\begin{lemma}[universal moment criteria inherit the margin]\label{lem:wh-T}
Let $\Phi:\R^k\to\R$ be such that $\Phi(m_1(B),\dots,m_k(B))<1$ implies $\|B\|<1$ for every nonnegative trace-class $B$
(a \emph{universal} criterion of order $k$). If $\beta_1=\|B_T\|<1$ and $v_1$ is a unit top eigenvector, then
$B':=B_T+(1-\beta_1)v_1\otimes v_1$ is nonnegative, trace class, with $\|B'\|=1$, so $\Phi(m(B'))\ge1$, and
\[
|m_j(B')-m_j(B_T)|=1-\beta_1^j\le j(1-\beta_1)\le\frac{j\,\lambda_{\mathrm A}(T)}{c_P}\qquad(1\le j\le k).
\]
Hence the point $m(B_T)$ lies within $k(1-\beta_1)$, in each coordinate, of the set $\{\Phi\ge1\}$: if $\Phi$ is $L$-Lipschitz
(sup-norm), then $1-\Phi(m(B_T))\le L\,k(1-\beta_1)\le Lk\lambda_{\mathrm A}(T)/c_P$.
\end{lemma}
\begin{proof}
$B'$ has the eigenvalues of $B_T$ with $\beta_1$ replaced by $1$; the rest is \eqref{eq:wh-sandwich} and $1-x^j\le j(1-x)$ on $[0,1]$.
\end{proof}

\begin{lemma}[any continuous universal criterion inherits the margin]\label{lem:wh-Tprime}
Let $\mathcal F$ be the set of forms $\mathfrak q=\fa|_{\PT}+\mathfrak b$ with $\mathfrak b$ a bounded symmetric form on $\PT$, and let
$\Phi:\mathcal F\to\R$ satisfy $|\Phi(\mathfrak q)-\Phi(\mathfrak q')|\le L\|\mathfrak b-\mathfrak b'\|$ and the universal property
``$\Phi(\mathfrak q)>0\Rightarrow\inf_{\|u\|=1}\mathfrak q(u)>0$''. If $\lambda_{\mathrm A}(T)\ge0$ then $\Phi(\fa|_{\PT})\le L\,\lambda_{\mathrm A}(T)$.
\end{lemma}
\begin{proof}
$\mathfrak q:=\fa|_{\PT}-\lambda_{\mathrm A}(T)\langle\cdot,h^*_T\rangle h^*_T\in\mathcal F$ has bottom $0$ (it is $\ge0$ and vanishes at $h^*_T$),
so $\Phi(\mathfrak q)\le0$, and $\|\mathfrak b\|=\lambda_{\mathrm A}(T)$.
\end{proof}

\begin{remark}\label{rem:wh-mom}
(1) The $m_k$ are computable from the symbol alone (through $\fp^{-1}$, i.e.\ the inverse of the compressed positive
part), and the criterion of Lemma~\ref{lem:wh-mom} is exact; but $k^*:=\min\{k:m_k<1\}$ is $20$, $\approx10^5$, $\approx10^{10}$
at $T=\frac12\log3,\log2,\frac12\log5$, with $m_k^{1/k}\in[1.06,6.7]$ for $k\le6$, and at $k^*$ the margin $1-m_{k^*}^{1/k^*}$ equals
$1-\beta_1$ to all digits (NUMERICAL, R186\_045 \S3.3). This is the determinant $\det(1-sB_T)=\exp(-\sum_ks^km_k/k)$ of
R186\_034 read in its power series. (2) Lemmas~\ref{lem:wh-T} and \ref{lem:wh-Tprime} concern criteria valid for all
forms (resp.\ all nonnegative trace-class $B$); a criterion valid only on symbols of the arithmetic class is not covered,
because $B'$ and $\mathfrak q$ need not come from such a symbol. No criterion of that kind is formulated anywhere in the
corpus (R186\_045 \S3.4).
\end{remark}

\subsection{Partition into old and new negativity, and contact thresholds}\label{subsec:wh-oldnew}

\begin{proposition}[old/new partition]\label{prop:wh-oldnew}
Let $T_0<T$, $U_0:=\{\sigma_{T_0}<0\}$,
\[
\mathrm{Neg}^{\mathrm{new}}_T:=R_T\,\sigma_-\mathbf 1_{U_T\setminus U_0}(D)R_T\ \ge0\quad\text{(bounded, trace class)}
\]
and $\mathfrak a^{(U_0)}_T:=\fa+\mathrm{Neg}^{\mathrm{new}}_T$, a closed form on $\DT$ (the form of the symbol whose only
negativity is that of $\sT$ on $U_0$). Then on $\PT\cap\DT$:
\emph{(i)} $\lambda_{\mathrm A}(T)\ge0\iff\mathfrak a^{(U_0)}_T\ge\mathrm{Neg}^{\mathrm{new}}_T$; in particular $\lambda_{\mathrm A}(T)\ge0$ implies
$\mathfrak a^{(U_0)}_T\ge0$ and $\lambda_{\mathrm A}(T)>0$ implies $\mathfrak a^{(U_0)}_T\ge\lambda_{\mathrm A}(T)$.
\emph{(ii)} If $\mathfrak a^{(U_0)}_T\ge c>0$ on $\PT$, the pencil top
\[
\beta_{\mathrm{new}}:=\sup_u\mathrm{Neg}^{\mathrm{new}}_T(u)/\mathfrak a^{(U_0)}_T(u)
\]
is finite
and attained, $\lambda_{\mathrm A}(T)>0\iff\beta_{\mathrm{new}}<1$, $\lambda_{\mathrm A}(T)\ge0\iff\beta_{\mathrm{new}}\le1$, and when
$\lambda_{\mathrm A}(T)\ge0$, with $u^*$ the unit pencil extremal,
\[
\lambda_{\mathrm A}(T)/\mathfrak a^{(U_0)}_T(u^*)\ \le\ 1-\beta_{\mathrm{new}}\ \le\ \lambda_{\mathrm A}(T)/\mathfrak a^{(U_0)}_T(h^*_T).
\]
Hence $\lambda_{\mathrm A}(T)>0\iff[\mathfrak a^{(U_0)}_T\ge c>0]\wedge[\beta_{\mathrm{new}}<1]$.
\end{proposition}
\begin{proof}
$\sT=\sT\mathbf 1_{\R\setminus(U_T\setminus U_0)}-\sigma_-\mathbf 1_{U_T\setminus U_0}$ is an identity of multipliers, so
$\fa=\mathfrak a^{(U_0)}_T-\mathrm{Neg}^{\mathrm{new}}_T$; the rest is Theorem~\ref{thm:wh-pencil} with $(\mathfrak a^{(U_0)}_T,\mathrm{Neg}^{\mathrm{new}}_T)$ in
place of $(\Pos,\Neg)$ (the closedness and lower bound of $\mathfrak a^{(U_0)}_T$ follow from those of $\fa$ and the boundedness of the added term).
\end{proof}
The first bracket has a margin of its own at the computed pairs ($\lambda_{\min}(\mathfrak a^{(U_0)}_T|\PT)=0.03$--$0.19$) and holds
equally for the modified symbols of R186\_114; the second inherits exactly: $1-\beta_{\mathrm{new}}=\lambda_{\mathrm A}(T)/\mathfrak a^{(U_0)}_T(h^*_T)$ to all
digits, and the ``new'' negativity seen by the old prolates is $O(1)$, not $o(1)$, dominated by growth and fission of components
adjacent to $U_0$ (NUMERICAL, R186\_049). The proposition is a reformulation; it is entered as structure only.

\begin{lemma}[a contact threshold does not change the form]\label{lem:wh-birth}
Let $n\ge2$ and $T_n:=\frac12\log n$. For every $F\in L^2(I_{T_n})$, $\frac1{2\pi}\int\cos(\tau\log n)|\Fh|^2d\tau=\Rel\cF(\log n)=0$.
Hence the form with symbol $\sigma_{T_n}-v_n\cos(\tau\log n)$ (the symbol of every window $T\in(T_n,T_{n'})$, $n'$ the next contact)
coincides with $\mathfrak a_{T_n}$ on $L^2(I_{T_n})$, and $\lambda_{\mathrm A}$, $\lamW$ are continuous at $T_n$ (Lemma~\ref{lem:contA}).
\end{lemma}
\begin{proof}
$\cF$ is continuous and supported in $[-2T_n,2T_n]$, and $\log n=2T_n$; the identity is Plancherel
($\widehat{\cF}=|\Fh|^2$). This is the computation in the proof of Lemma~\ref{lem:monoT}.
\end{proof}

\begin{lemma}[newborn components have Shannon number below $\frac12$]\label{lem:wh-shannon}
Let $T_n=\frac12\log n$, $\sigma^-:=\sigma_{T_n}$ (without $n$) and $\sigma^+:=\sigma^--v_n\cos(\tau\log n)$ (with $n$). If $J\subset(0,\infty)$
is a connected component of $\{\sigma^+<0\}$ disjoint from $\{\sigma^-<0\}$, then $\cos(\tau\log n)>0$ on $J$, $|J|<\pi/\log n$, and
$T_n|J|/\pi<\frac12$ (equivalently $\tr S_{J\cup(-J)}<1$ at the window $T_n$).
\end{lemma}
\begin{proof}
On $J$, $\sigma^-\ge0$ and $v_n\cos(\tau\log n)>\sigma^-\ge0$, so the connected $J$ lies in one half-period $\{\cos(\tau\log n)>0\}$,
of length $\pi/\log n$; then $T_n|J|/\pi<T_n/\log n=\frac12$.
\end{proof}
Such components are born with mass $O(1)$ (not small) and invisible at the instant of birth (Lemma~\ref{lem:wh-birth}); afterwards
the cost of the new contact is $v_n\Rel c_{h^*}(\log n)$ at first order, a quantity of the minimizer, and it exhausts the margin
$\lambda_{\mathrm A}(T_n)$ within a fraction of the interval to the next contact (NUMERICAL, R186\_048 \S4). No perturbative
continuation argument with a margin of its own follows from these two lemmas (R186\_048 Prop.~4.3: the margin is $\lambda_{\mathrm A}(T_n)$).

\subsection{The negative index of the truncated form by parity sectors}\label{subsec:wh-index}

We use the notation of Theorem~\ref{thm:parityschur}: $\varphi_{\mathrm{ev}}=\cosh\frac t2\mathbf 1_{\IT}$,
$\varphi_{\mathrm{odd}}=\sinh\frac t2\mathbf 1_{\IT}$, $\ell_{\mathrm{par}}=\langle\cdot,\varphi_{\mathrm{par}}\rangle$,
$\varepsilon_{\mathrm{ev}}=+1$, $\varepsilon_{\mathrm{odd}}=-1$, $\mathcal Q^{\mathrm{par}}_T=\ker\ell_{\mathrm{par}}\cap\DT^{\mathrm{par}}$,
$h_{\mathrm{par}}$ and $\mathfrak g_{\mathrm{par}}=\fa(h_{\mathrm{par}})+2\varepsilon_{\mathrm{par}}$ (Proposition~\ref{prop:polarcap}).
By Theorem~\ref{thm:parityschur}(o), on real $F=F_{\mathrm e}+F_{\mathrm o}\in\DT$,
\begin{equation}\label{eq:wh-aW}
\fa(F)=W_T(F)-2\ell_{\mathrm{ev}}(F_{\mathrm e})^2+2\ell_{\mathrm{odd}}(F_{\mathrm o})^2 .
\end{equation}

\begin{proposition}[index count by sectors]\label{prop:wh-index}
\emph{(i)} If $0$ is not in the spectrum of the operator of $\fa|\mathcal Q^{\mathrm{par}}_T$ in both sectors (in particular if $\lambda_{\mathrm A}(T)>0$), then
\[
n_-(\fa|\DT)=n_-(\fa|\mathcal Q^{\mathrm{ev}}_T)+[\mathfrak g_{\mathrm{ev}}(T)<2]+n_-(\fa|\mathcal Q^{\mathrm{odd}}_T)+[\mathfrak g_{\mathrm{odd}}(T)<-2].
\]
\emph{(ii)} For $0<T\le\frac12\log5$: $n_-(\fa|\DT^{\mathrm{odd}})=0$, $n_-(\fa|\DT^{\mathrm{ev}})\le1$, and $\fa\ge10^{-18}\|\cdot\|^2$ on
$\{F\in\DT:\ell_{\mathrm{ev}}(F)=0\}$, a space of codimension one in $\DT$ containing $\PT$. If moreover $\fa(\varphi_{\mathrm{ev}})<0$, then
$n_-(\fa|\DT)=1$, the negative direction is even, and $\fa(h_{\mathrm{ev}})<0$, i.e.\ $\mathfrak g_{\mathrm{ev}}(T)<2$.
\emph{(iii)} $\fa(\varphi_{\mathrm{ev}})\le\Gamma_T(\varphi_{\mathrm{ev}})$ for every $T$, where
$\Gamma_T(\varphi_{\mathrm{ev}})=C_0\|\varphi_{\mathrm{ev}}\|^2+\int_0^{2T}\big(\|\varphi_{\mathrm{ev}}\|^2-c_{\varphi_{\mathrm{ev}}}(u)\big)K(u)du+\|\varphi_{\mathrm{ev}}\|^2\int_{2T}^\infty K$
is an explicit function of $T$ alone.
\end{proposition}
\begin{proof}
(i) Theorem~\ref{thm:parityschur}(a1): $\DT^{\mathrm{par}}=\mathcal Q^{\mathrm{par}}_T\oplus\R h_{\mathrm{par}}$, $\fa$-orthogonally, so
$\fa(G+xh)=\fa(G)+x^2\fa(h_{\mathrm{par}})$ and $n_-(\fa|\DT^{\mathrm{par}})=n_-(\fa|\mathcal Q^{\mathrm{par}}_T)+[\fa(h_{\mathrm{par}})<0]$; the sectors are
$\fa$-orthogonal (Theorem~\ref{thm:parityschur}(o)), and $\fa(h_{\mathrm{par}})=\mathfrak g_{\mathrm{par}}-2\varepsilon_{\mathrm{par}}$.
(ii) By Theorem~\ref{thm:log5W} and Lemma~\ref{lem:monoT}, $W_T\ge10^{-18}\|\cdot\|^2$ on $\DT$ for $T\le\frac12\log5$. By \eqref{eq:wh-aW},
on odd $F$, $\fa(F)=W_T(F)+2\ell_{\mathrm{odd}}(F)^2\ge0$; on even $F$ with $\ell_{\mathrm{ev}}(F)=0$, $\fa(F)=W_T(F)$; so on $\ker\ell_{\mathrm{ev}}$,
$\fa(F)=W_T(F_{\mathrm e})+W_T(F_{\mathrm o})+2\ell_{\mathrm{odd}}(F_{\mathrm o})^2\ge10^{-18}\|F\|^2$, and $n_-(\fa|\DT^{\mathrm{ev}})\le1$ by min--max
($\ker\ell_{\mathrm{ev}}$ has codimension $1$ in $\DT^{\mathrm{ev}}$). If $\fa(\varphi_{\mathrm{ev}})<0$ then $n_-\ge1$ in the even sector, hence $=1$,
and $n_-(\fa|\mathcal Q^{\mathrm{ev}}_T)=0$ (as $\mathcal Q^{\mathrm{ev}}_T\subset\ker\ell_{\mathrm{ev}}$) forces $\fa(h_{\mathrm{ev}})<0$ by (i).
(iii) $\fa=\Gamma_T+\PiT$ with $\PiT(F)=-\sum v_n\Rel\cF(\log n)\le0$ for $F\ge0$ ($\cF\ge0$ then), and \eqref{eq:Aform} with $\NT=\emptyset$
gives the formula ($\Rel c_{\varphi_{\mathrm{ev}}}=c_{\varphi_{\mathrm{ev}}}$, $c_{\varphi_{\mathrm{ev}}}=0$ beyond $2T$).
\end{proof}

\begin{remark}\label{rem:wh-index}
(1) $\Gamma_T(\varphi_{\mathrm{ev}})<0$ for $T\ge0.12$ ($-0.047$ at $0.12$, $-0.84$ at $0.25$, $-3.31$ at $\frac12\log5$) is a one-dimensional
evaluation of the explicit expression in (iii); it is NUMERICAL (certifiable by interval arithmetic, not done; R186\_028 Cor.~1.3).
With it, $n_-(\fa|\DT)=1$ on $[0.12,\frac12\log5]$ with even negative direction (R186\_028 Thm 1.4; hostile R186\_106/107: the upper bound
is a two-line consequence of \eqref{eq:wh-aW} and $\lamW\ge0$, and its range is exactly that of Theorem~\ref{thm:log5W}, not of \cite{Zhu2026}).
(2) \emph{Not a lever.} For $T\ge0.12$ with both sectors nondegenerate, (i) gives $[n_-(\fa|\DT)=1\wedge\fa(h_{\mathrm{ev}})<0]\iff
[\lambda_{\mathrm A}(T)>0\wedge\mathfrak g_{\mathrm{odd}}(T)\ge-2]$ (R186\_028 Cor.~1.3(iii)); the index statement contains the positivity
statement and nothing weaker: without the sign of $\fa(h_{\mathrm{ev}})$, $n_-=1$ is compatible with $n_-(\fa|\mathcal Q^{\mathrm{ev}}_T)=1$.
(3) \emph{The polar resolvent.} If $\fa$ is invertible on $\DT^{\mathrm{par}}$ and $\sigma_{\mathrm{par}}:=\langle\fa^{-1}\varphi_{\mathrm{par}},\varphi_{\mathrm{par}}\rangle\ne0$,
then $h_{\mathrm{par}}=\fa^{-1}\varphi_{\mathrm{par}}/\sigma_{\mathrm{par}}$, $\fa(h_{\mathrm{par}})=1/\sigma_{\mathrm{par}}$ and
$\sigma_{\mathrm{par}}=1/(\mathfrak g_{\mathrm{par}}-2\varepsilon_{\mathrm{par}})$; so the $2\times2$ matrix $\big(\langle\fa^{-1}\ell_i,\ell_j\rangle\big)$ of
the Tate functionals $(A,B)$ has spectrum $\{2/(\mathfrak g_{\mathrm{ev}}-2),\,2/(\mathfrak g_{\mathrm{odd}}+2)\}$ (R186\_028 Thm 2.1, algebra from
Proposition~\ref{prop:polarcap}). The eigenvalue $\approx-1$ observed in R186\_026 is $-1-\mathfrak g_{\mathrm{ev}}/2+O(\mathfrak g_{\mathrm{ev}}^2)$,
with $\mathfrak g_{\mathrm{ev}}\asymp\lamW\|h_{\mathrm{ev}}\|^2$ (Theorem~\ref{thm:parityschur}(a4)); it is the inverse of the negative square
of the pole (Lemma~\ref{lem:polar}), explained and closed.
\end{remark}

\subsection{Labels, and what is not asserted}\label{subsec:wh-labels}

\begin{remark}[labels]\label{rem:wh-labels}
Proved: Lemma~\ref{lem:wh-cP}, Theorem~\ref{thm:wh-pencil}, Proposition~\ref{prop:wh-basis}, Theorem~\ref{thm:wh-angles},
Corollary~\ref{cor:wh-nogap}, Proposition~\ref{prop:wh-notangles}, Lemma~\ref{lem:wh-krein}, Corollary~\ref{cor:wh-step},
Lemma~\ref{lem:wh-W}, Theorem~\ref{thm:wh-necessary}, Lemmas~\ref{lem:wh-neg}, \ref{lem:wh-pos}, Theorem~\ref{thm:wh-LW} (as an
inequality with its constants defined), Theorem~\ref{thm:wh-schur} (conditional only on $A_{WW}\ge c>0$), Lemma~\ref{lem:wh-Q},
Proposition~\ref{prop:wh-K}, Lemmas~\ref{lem:wh-mom}, \ref{lem:wh-T}, \ref{lem:wh-Tprime}, Proposition~\ref{prop:wh-oldnew},
Lemmas~\ref{lem:wh-birth}, \ref{lem:wh-shannon}, Proposition~\ref{prop:wh-index}. Numerical and not proved: every value quoted in
Remarks~\ref{rem:wh-pencil}(2)(3), \ref{rem:wh-model}, \ref{rem:wh-LWnum}, \ref{rem:wh-mom}, \ref{rem:wh-index}(1) and after
Propositions~\ref{prop:wh-basis}, \ref{prop:wh-oldnew}; the certified constants of Remark~\ref{rem:wh-LWnum}(1) are proved modulo
convergent floating-point quadrature. Argued only: the slope law of Remark~\ref{rem:wh-model}(1) and the first-order law after
Proposition~\ref{prop:wh-K}. Cited and not used in any proof: the asymptotics of Slepian, Fuchs, Widom and Landau--Widom.
\end{remark}

\begin{remark}[what this section does not assert]\label{rem:wh-honest}
No statement above bounds $\lambda_{\mathrm A}(T)$, $1-\beta_1$, $1-\rho_T$ or $1-\beta_{\mathrm{new}}$ from below at any window beyond
Theorem~\ref{thm:log5A}, and none asserts a sign for any $T$. Each exact reduction (pencil, Schur complement over the prolate block,
old/new partition, moments) replaces $\lambda_{\mathrm A}(T)$ by a quantity that is $\lambda_{\mathrm A}(T)$ divided by a number of size
$O(1)$ (\eqref{eq:wh-sandwich}, Theorem~\ref{thm:wh-schur}(iii), Proposition~\ref{prop:wh-oldnew}(ii), Lemma~\ref{lem:wh-T}); the
quantities with a margin of their own ($c_W$, $\lambda_{\min}(H_T)$, $\lambda_{\min}(\mathfrak a^{(U_0)}_T)$, $c_P$) are the same for $\sT$
and for the modified symbols of R186\_114 at the scale where those change sign (NUMERICAL), as Lemmas~\ref{lem:wh-T} and
\ref{lem:wh-Tprime} say they must be for any universal criterion. What is new relative to Proposition~\ref{prop:negtrace} is the
structure: the sign is the position of the last inertia crossing of $\beta\mapsto R_T(\sigma_+-\sigma_-/\beta)(D)R_T$ on $\PT$ relative to
$\beta=1$; it is not a function of $U_T$, of $\sigma_-$, of the angles of $(R_T,\mathbf 1_{U_T}(D))$, of any finite list of moments, or of
the index by sectors.
\end{remark}

\ifdefined\DraftStandalone\end{document}\fi
